% Options for packages loaded elsewhere
\PassOptionsToPackage{unicode}{hyperref}
\PassOptionsToPackage{hyphens}{url}
%
\documentclass[
]{article}
\usepackage{amsmath,amssymb}
\usepackage{lmodern}
\usepackage{iftex}
\ifPDFTeX
  \usepackage[T1]{fontenc}
  \usepackage[utf8]{inputenc}
  \usepackage{textcomp} % provide euro and other symbols
\else % if luatex or xetex
  \usepackage{unicode-math}
  \defaultfontfeatures{Scale=MatchLowercase}
  \defaultfontfeatures[\rmfamily]{Ligatures=TeX,Scale=1}
\fi
% Use upquote if available, for straight quotes in verbatim environments
\IfFileExists{upquote.sty}{\usepackage{upquote}}{}
\IfFileExists{microtype.sty}{% use microtype if available
  \usepackage[]{microtype}
  \UseMicrotypeSet[protrusion]{basicmath} % disable protrusion for tt fonts
}{}
\makeatletter
\@ifundefined{KOMAClassName}{% if non-KOMA class
  \IfFileExists{parskip.sty}{%
    \usepackage{parskip}
  }{% else
    \setlength{\parindent}{0pt}
    \setlength{\parskip}{6pt plus 2pt minus 1pt}}
}{% if KOMA class
  \KOMAoptions{parskip=half}}
\makeatother
\usepackage{xcolor}
\IfFileExists{xurl.sty}{\usepackage{xurl}}{} % add URL line breaks if available
\IfFileExists{bookmark.sty}{\usepackage{bookmark}}{\usepackage{hyperref}}
\hypersetup{
  pdftitle={Ejercicio 7},
  pdfauthor={Alex Pacheco - Francisco Salfate - Xavier Munoz - Elena Guzman},
  hidelinks,
  pdfcreator={LaTeX via pandoc}}
\urlstyle{same} % disable monospaced font for URLs
\usepackage[margin=1in]{geometry}
\usepackage{color}
\usepackage{fancyvrb}
\newcommand{\VerbBar}{|}
\newcommand{\VERB}{\Verb[commandchars=\\\{\}]}
\DefineVerbatimEnvironment{Highlighting}{Verbatim}{commandchars=\\\{\}}
% Add ',fontsize=\small' for more characters per line
\usepackage{framed}
\definecolor{shadecolor}{RGB}{248,248,248}
\newenvironment{Shaded}{\begin{snugshade}}{\end{snugshade}}
\newcommand{\AlertTok}[1]{\textcolor[rgb]{0.94,0.16,0.16}{#1}}
\newcommand{\AnnotationTok}[1]{\textcolor[rgb]{0.56,0.35,0.01}{\textbf{\textit{#1}}}}
\newcommand{\AttributeTok}[1]{\textcolor[rgb]{0.77,0.63,0.00}{#1}}
\newcommand{\BaseNTok}[1]{\textcolor[rgb]{0.00,0.00,0.81}{#1}}
\newcommand{\BuiltInTok}[1]{#1}
\newcommand{\CharTok}[1]{\textcolor[rgb]{0.31,0.60,0.02}{#1}}
\newcommand{\CommentTok}[1]{\textcolor[rgb]{0.56,0.35,0.01}{\textit{#1}}}
\newcommand{\CommentVarTok}[1]{\textcolor[rgb]{0.56,0.35,0.01}{\textbf{\textit{#1}}}}
\newcommand{\ConstantTok}[1]{\textcolor[rgb]{0.00,0.00,0.00}{#1}}
\newcommand{\ControlFlowTok}[1]{\textcolor[rgb]{0.13,0.29,0.53}{\textbf{#1}}}
\newcommand{\DataTypeTok}[1]{\textcolor[rgb]{0.13,0.29,0.53}{#1}}
\newcommand{\DecValTok}[1]{\textcolor[rgb]{0.00,0.00,0.81}{#1}}
\newcommand{\DocumentationTok}[1]{\textcolor[rgb]{0.56,0.35,0.01}{\textbf{\textit{#1}}}}
\newcommand{\ErrorTok}[1]{\textcolor[rgb]{0.64,0.00,0.00}{\textbf{#1}}}
\newcommand{\ExtensionTok}[1]{#1}
\newcommand{\FloatTok}[1]{\textcolor[rgb]{0.00,0.00,0.81}{#1}}
\newcommand{\FunctionTok}[1]{\textcolor[rgb]{0.00,0.00,0.00}{#1}}
\newcommand{\ImportTok}[1]{#1}
\newcommand{\InformationTok}[1]{\textcolor[rgb]{0.56,0.35,0.01}{\textbf{\textit{#1}}}}
\newcommand{\KeywordTok}[1]{\textcolor[rgb]{0.13,0.29,0.53}{\textbf{#1}}}
\newcommand{\NormalTok}[1]{#1}
\newcommand{\OperatorTok}[1]{\textcolor[rgb]{0.81,0.36,0.00}{\textbf{#1}}}
\newcommand{\OtherTok}[1]{\textcolor[rgb]{0.56,0.35,0.01}{#1}}
\newcommand{\PreprocessorTok}[1]{\textcolor[rgb]{0.56,0.35,0.01}{\textit{#1}}}
\newcommand{\RegionMarkerTok}[1]{#1}
\newcommand{\SpecialCharTok}[1]{\textcolor[rgb]{0.00,0.00,0.00}{#1}}
\newcommand{\SpecialStringTok}[1]{\textcolor[rgb]{0.31,0.60,0.02}{#1}}
\newcommand{\StringTok}[1]{\textcolor[rgb]{0.31,0.60,0.02}{#1}}
\newcommand{\VariableTok}[1]{\textcolor[rgb]{0.00,0.00,0.00}{#1}}
\newcommand{\VerbatimStringTok}[1]{\textcolor[rgb]{0.31,0.60,0.02}{#1}}
\newcommand{\WarningTok}[1]{\textcolor[rgb]{0.56,0.35,0.01}{\textbf{\textit{#1}}}}
\usepackage{graphicx}
\makeatletter
\def\maxwidth{\ifdim\Gin@nat@width>\linewidth\linewidth\else\Gin@nat@width\fi}
\def\maxheight{\ifdim\Gin@nat@height>\textheight\textheight\else\Gin@nat@height\fi}
\makeatother
% Scale images if necessary, so that they will not overflow the page
% margins by default, and it is still possible to overwrite the defaults
% using explicit options in \includegraphics[width, height, ...]{}
\setkeys{Gin}{width=\maxwidth,height=\maxheight,keepaspectratio}
% Set default figure placement to htbp
\makeatletter
\def\fps@figure{htbp}
\makeatother
\setlength{\emergencystretch}{3em} % prevent overfull lines
\providecommand{\tightlist}{%
  \setlength{\itemsep}{0pt}\setlength{\parskip}{0pt}}
\setcounter{secnumdepth}{-\maxdimen} % remove section numbering
\ifLuaTeX
  \usepackage{selnolig}  % disable illegal ligatures
\fi

\title{Ejercicio 7}
\author{Alex Pacheco - Francisco Salfate - Xavier Munoz - Elena Guzman}
\date{2022-10-04}

\begin{document}
\maketitle

\begin{enumerate}
\def\labelenumi{\arabic{enumi}.}
\tightlist
\item
  Observando los datos, la memorista sospecha que hay diferencias
  significativas en el tiempo de ejecución entre las versiones A y B del
  algoritmo cuando las instancias tienen 70 o más nodos. ¿Los datos
  respaldan la intuición de la memorista? Para responder, filtren los
  datos para tener las instancias con 70 o más nodos y seleccionen las
  columnas de los tiempos de ejecución de las versiones A y B en formato
  ancho. Usando como semilla el valor 73, obtengan muestras aleatorias
  independientes de 24 tiempos registrados por la versión A y 20 tiempos
  registrados por la versión B del algoritmo. Lleven los datos a formato
  largo y utilicen la escalera de potencias de Tukey para analizar las
  muestras obtenidas.
\end{enumerate}

\hypertarget{hipuxf3tesis}{%
\subsection{Hipótesis:}\label{hipuxf3tesis}}

\hypertarget{h0-los-tiempos-de-ejecuciuxf3n-entre-la-versiuxf3n-a-y-b-del-algoritmo-cuando-estas-poseen-instancias-de-70-o-muxe1s-nodos-no-poseen-diferencias-significativas}{%
\subsubsection{\texorpdfstring{\textbf{H0 : Los tiempos de ejecución
entre la versión A y B del algoritmo cuando estas poseen instancias de
70 o más nodos, no poseen diferencias
significativas}}{H0 : Los tiempos de ejecución entre la versión A y B del algoritmo cuando estas poseen instancias de 70 o más nodos, no poseen diferencias significativas}}\label{h0-los-tiempos-de-ejecuciuxf3n-entre-la-versiuxf3n-a-y-b-del-algoritmo-cuando-estas-poseen-instancias-de-70-o-muxe1s-nodos-no-poseen-diferencias-significativas}}

\[ \mu_{dif} = 0 \]

\hypertarget{ha-los-tiempos-de-ejecuciuxf3n-entre-la-versiuxf3n-a-y-b-del-algoritmo-cuando-estas-poseen-instancias-de-70-o-muxe1s-nodos-poseen-diferencias-significativas.}{%
\subsubsection{\texorpdfstring{\textbf{Ha : Los tiempos de ejecución
entre la versión A y B del algoritmo cuando estas poseen instancias de
70 o más nodos, poseen diferencias
significativas.}}{Ha : Los tiempos de ejecución entre la versión A y B del algoritmo cuando estas poseen instancias de 70 o más nodos, poseen diferencias significativas.}}\label{ha-los-tiempos-de-ejecuciuxf3n-entre-la-versiuxf3n-a-y-b-del-algoritmo-cuando-estas-poseen-instancias-de-70-o-muxe1s-nodos-poseen-diferencias-significativas.}}

\[ \mu_{dif}  \neq 0 \]

\begin{Shaded}
\begin{Highlighting}[]
\FunctionTok{set.seed}\NormalTok{(}\DecValTok{73}\NormalTok{)}
\NormalTok{datos }\OtherTok{\textless{}{-}}\FunctionTok{read.csv}\NormalTok{(}\AttributeTok{file=}\StringTok{"EP07 Datos.csv"}\NormalTok{)}
\NormalTok{over70Nodos }\OtherTok{=}\NormalTok{ datos }\SpecialCharTok{\%\textgreater{}\%} \FunctionTok{filter}\NormalTok{ (n.nodos }\SpecialCharTok{\textgreater{}=} \DecValTok{70}\NormalTok{)}
\end{Highlighting}
\end{Shaded}

\hypertarget{selecciuxf3n-de-muestras-para-el-algoritmo-a-y-b}{%
\subsubsection{Selección de muestras para el algoritmo A y
B}\label{selecciuxf3n-de-muestras-para-el-algoritmo-a-y-b}}

Primero, se seleccionan los datos que posean 70 nodos o más, los cuales
están en formato ancho. Luego, se seleccionan de forma aleatoria 24
muestras para el tiempo del algoritmo A y 20 para el algoritmo B.

\begin{Shaded}
\begin{Highlighting}[]
\CommentTok{\# Data frame con n.nodos \textgreater{}= 70 }
\NormalTok{muestra }\OtherTok{\textless{}{-}}\NormalTok{ datos[}\FunctionTok{sample}\NormalTok{(}\FunctionTok{which}\NormalTok{(datos}\SpecialCharTok{$}\NormalTok{n.nodos }\SpecialCharTok{\textgreater{}=} \DecValTok{70}\NormalTok{), }\DecValTok{44}\NormalTok{), ]}
\NormalTok{muestra\_t\_A }\OtherTok{\textless{}{-}}\NormalTok{ muestra[, }\DecValTok{4}\NormalTok{][}\DecValTok{1}\SpecialCharTok{:}\DecValTok{24}\NormalTok{]}
\NormalTok{muestra\_t\_B }\OtherTok{\textless{}{-}}\NormalTok{ muestra[, }\DecValTok{5}\NormalTok{][}\DecValTok{25}\SpecialCharTok{:}\DecValTok{44}\NormalTok{]}
\end{Highlighting}
\end{Shaded}

\hypertarget{prueba-de-normalidad-para-cada-muestra}{%
\subsubsection{Prueba de normalidad para cada
muestra}\label{prueba-de-normalidad-para-cada-muestra}}

\begin{Shaded}
\begin{Highlighting}[]
\NormalTok{shapiroMuestraA }\OtherTok{\textless{}{-}} \FunctionTok{shapiro.test}\NormalTok{(muestra\_t\_A)}
\NormalTok{shapiroMuestraB }\OtherTok{\textless{}{-}} \FunctionTok{shapiro.test}\NormalTok{(muestra\_t\_B)}
\FunctionTok{c}\NormalTok{(shapiroMuestraA}\SpecialCharTok{$}\NormalTok{p.value, shapiroMuestraB}\SpecialCharTok{$}\NormalTok{p.value)}
\end{Highlighting}
\end{Shaded}

\begin{verbatim}
## [1] 0.001783196 0.033611703
\end{verbatim}

Considerando un nivel de significancia de alpha = 0.05, al aplicar
shapiro.test a las muestras del algoritmo A y del algoritmo B, se da
cuenta que el valor p obtenido en ambas pruebas, es menor al alpha. En
consecuencia, dichas muestras no siguen una distribución cercana a la
normal y es necesario realizar una transformación de los datos para
poder analizarlos. En este caso, se utilizará la escalera de potencias
de Tukey. Sin embargo, los datos se deben llevar a formato largo para
poder utilizar la escalera de potencias de Tukey para analizar las
muestras obtenidas.

\begin{Shaded}
\begin{Highlighting}[]
\NormalTok{id }\OtherTok{\textless{}{-}} \FunctionTok{factor}\NormalTok{(}\DecValTok{1}\SpecialCharTok{:}\DecValTok{44}\NormalTok{)}
\NormalTok{algoritmos }\OtherTok{\textless{}{-}} \FunctionTok{c}\NormalTok{(}
                \FunctionTok{replicate}\NormalTok{(}\DecValTok{24}\NormalTok{, }\StringTok{"A"}\NormalTok{),}
                \FunctionTok{replicate}\NormalTok{(}\DecValTok{20}\NormalTok{, }\StringTok{"B"}\NormalTok{)}
\NormalTok{                )}
\NormalTok{tiempos }\OtherTok{\textless{}{-}} \FunctionTok{c}\NormalTok{(muestra\_t\_A, muestra\_t\_B)}
\NormalTok{data }\OtherTok{\textless{}{-}} \FunctionTok{data.frame}\NormalTok{(}
\NormalTok{                    id }\OtherTok{\textless{}{-}}\NormalTok{ id,}
\NormalTok{                    algoritmos }\OtherTok{\textless{}{-}}\NormalTok{ algoritmos,}
\NormalTok{                    tiempos }\OtherTok{\textless{}{-}}\NormalTok{ tiempos}
\NormalTok{                    )}
\end{Highlighting}
\end{Shaded}

\hypertarget{transformada-de-tukey}{%
\subsubsection{Transformada de tukey}\label{transformada-de-tukey}}

\begin{Shaded}
\begin{Highlighting}[]
\NormalTok{transformada }\OtherTok{\textless{}{-}} \FunctionTok{transformTukey}\NormalTok{(}
\NormalTok{                            data}\SpecialCharTok{$}\NormalTok{tiempos,}
\NormalTok{                            start }\OtherTok{\textless{}{-}} \SpecialCharTok{{-}}\DecValTok{2}\NormalTok{,}
\NormalTok{                            end }\OtherTok{\textless{}{-}} \DecValTok{2}\NormalTok{,}
\NormalTok{                            int }\OtherTok{\textless{}{-}} \FloatTok{0.001}\NormalTok{,}
\NormalTok{                            returnLambda  }\OtherTok{\textless{}{-}} \ConstantTok{FALSE}
\NormalTok{                            )}
\end{Highlighting}
\end{Shaded}

\begin{verbatim}
## 
##      lambda     W Shapiro.p.value
## 2524  0.523 0.964          0.1826
## 
## if (lambda >  0){TRANS = x ^ lambda} 
## if (lambda == 0){TRANS = log(x)} 
## if (lambda <  0){TRANS = -1 * x ^ lambda}
\end{verbatim}

Al observar el gráfico Q-Q obtenido al aplicar la escalera de potencias
de Tukey, se puede observar que los datos se ajustan a una distribución
normal. En consecuencia procede a aplicar una prueba t para dos muestras
no independientes.

Luego, se procede a verificar si se cumplen las condiciones necesarias
para poder aplicar la prueba t para dos muestras no independientes:

\begin{itemize}
\tightlist
\item
  Como las muestras obtenidas, tanto para el algoritmo A como para el
  algoritmo B, fueron escogidas aleatoriamente y son independientes
  entre si
\item
  Según la transformada de Tukey aplicada se tiene que los nuevos datos
  siguen una distribución normal, esto debido a que la prueba de shapiro
  da un valor 0.1826 mayor al 0.05, que es nuestro valor de alfa.
\end{itemize}

Dado que se cumplen, se procede a aplicar la prueba ya mencionada:

\hypertarget{aplicaciuxf3n-de-t-test-dos-muestras-independientes}{%
\subsubsection{Aplicación de t test dos muestras
independientes}\label{aplicaciuxf3n-de-t-test-dos-muestras-independientes}}

\begin{Shaded}
\begin{Highlighting}[]
\NormalTok{data}\SpecialCharTok{$}\NormalTok{tiempos }\OtherTok{\textless{}{-}}\NormalTok{ transformada}
\NormalTok{alfa }\OtherTok{\textless{}{-}} \FloatTok{0.05}
\NormalTok{muestraATransformada }\OtherTok{\textless{}{-}}\NormalTok{ (data }\SpecialCharTok{\%\textgreater{}\%} \FunctionTok{filter}\NormalTok{(algoritmos....algoritmos }\SpecialCharTok{==} \StringTok{"A"}\NormalTok{))}\SpecialCharTok{$}\NormalTok{tiempos}
\NormalTok{muestraBTransformada }\OtherTok{\textless{}{-}}\NormalTok{ (data }\SpecialCharTok{\%\textgreater{}\%} \FunctionTok{filter}\NormalTok{(algoritmos....algoritmos }\SpecialCharTok{==} \StringTok{"B"}\NormalTok{))}\SpecialCharTok{$}\NormalTok{tiempos}
\FunctionTok{t.test}\NormalTok{(}
\NormalTok{        x }\OtherTok{\textless{}{-}}\NormalTok{ muestraATransformada,}
\NormalTok{        y }\OtherTok{\textless{}{-}}\NormalTok{ muestraBTransformada,}
        \AttributeTok{paired =} \ConstantTok{FALSE}\NormalTok{,}
\NormalTok{        alternative }\OtherTok{\textless{}{-}} \StringTok{"two.sided"}\NormalTok{,}
        \AttributeTok{conf.level =} \DecValTok{1} \SpecialCharTok{{-}}\NormalTok{ alfa}
\NormalTok{        )}\SpecialCharTok{$}\NormalTok{p.value}
\end{Highlighting}
\end{Shaded}

\begin{verbatim}
## [1] 0.02178821
\end{verbatim}

Dado que el valor p es menor a 0.025 (que es nuestro alfa/2 gracias a
que es de dos colas) entonces se rechaza la hipótesis nula en favor de
la hipótesis alternativa, por lo tanto, las sospechas de la memorista
son ciertas, en otras palabras, las diferencias entre los algoritmos A y
B para instancias de 70 o más nodos presentan diferencias
significativas.

\begin{enumerate}
\def\labelenumi{\arabic{enumi}.}
\setcounter{enumi}{1}
\tightlist
\item
  La memorista también sospecha que, al comparar las mismas instancias,
  las mejores soluciones encontradas por las versiones B y C tienen
  rendimientos distintos. ¿Estará en lo cierto? Para responder, filtren
  los datos para tener las instancias con 70 o más nodos y seleccionen
  las columnas con el mejor rendimiento de las versiones B y C en
  formato ancho. Usando como semilla el valor 71, obtengan una muestra
  aleatoria de 24 instancias. Lleven los datos a formato largo y
  utilicen una prueba no paramétrica apropiada para analizar las
  muestras obtenidas.
\end{enumerate}

\hypertarget{hipuxf3tesis-1}{%
\subsection{Hipótesis:}\label{hipuxf3tesis-1}}

\hypertarget{h0-las-mejores-soluciones-encontradas-para-la-versiuxf3n-b-y-c-del-algoritmo-cuando-estas-poseen-instancias-de-70-o-muxe1s-nodos-no-poseen-rendimientos-distintos}{%
\subsubsection{\texorpdfstring{\textbf{H0 : Las mejores soluciones
encontradas para la versión B y C del algoritmo cuando estas poseen
instancias de 70 o más nodos, no poseen rendimientos
distintos}}{H0 : Las mejores soluciones encontradas para la versión B y C del algoritmo cuando estas poseen instancias de 70 o más nodos, no poseen rendimientos distintos}}\label{h0-las-mejores-soluciones-encontradas-para-la-versiuxf3n-b-y-c-del-algoritmo-cuando-estas-poseen-instancias-de-70-o-muxe1s-nodos-no-poseen-rendimientos-distintos}}

\[ \mu_B  = \mu_C \]

\hypertarget{ha-las-mejores-soluciones-encontradas-para-la-versiuxf3n-b-y-c-del-algoritmo-cuando-estas-poseen-instancias-de-70-o-muxe1s-nodos-poseen-rendimientos-distintos}{%
\subsubsection{\texorpdfstring{\textbf{Ha : Las mejores soluciones
encontradas para la versión B y C del algoritmo cuando estas poseen
instancias de 70 o más nodos, poseen rendimientos
distintos}}{Ha : Las mejores soluciones encontradas para la versión B y C del algoritmo cuando estas poseen instancias de 70 o más nodos, poseen rendimientos distintos}}\label{ha-las-mejores-soluciones-encontradas-para-la-versiuxf3n-b-y-c-del-algoritmo-cuando-estas-poseen-instancias-de-70-o-muxe1s-nodos-poseen-rendimientos-distintos}}

\[ \mu_B \neq \mu_C \]

Dado que se quiere determinar si para una misma instancia existen
distintos rendimientos, se determina que se debe utilizar la prueba de
rangos con signos de Wilcoxon. Sin embargo, primero se debe determinar
si se cumplen las condiciones para poder aplicar dicha prueba.

\hypertarget{prueba-de-rangos-con-signos-de-wilcoxon}{%
\subsubsection{Prueba de rangos con signos de
Wilcoxon}\label{prueba-de-rangos-con-signos-de-wilcoxon}}

Primero se seleccionan 24 datos de forma aleatoria para luego aplicar la
prueba de rangos con signos de Wilcoxon

\begin{Shaded}
\begin{Highlighting}[]
\FunctionTok{set.seed}\NormalTok{(}\DecValTok{71}\NormalTok{)}
\NormalTok{muestra }\OtherTok{\textless{}{-}}\NormalTok{ datos[}\FunctionTok{sample}\NormalTok{(}\FunctionTok{which}\NormalTok{(datos}\SpecialCharTok{$}\NormalTok{n.nodos }\SpecialCharTok{\textgreater{}=} \DecValTok{70}\NormalTok{), }\DecValTok{24}\NormalTok{), ]}
\NormalTok{mejoresB }\OtherTok{\textless{}{-}}\NormalTok{ muestra}\SpecialCharTok{$}\NormalTok{mejor.B}
\NormalTok{mejoresC }\OtherTok{\textless{}{-}}\NormalTok{ muestra}\SpecialCharTok{$}\NormalTok{mejor.C}
\NormalTok{prueba }\OtherTok{\textless{}{-}} \FunctionTok{wilcox.test}\NormalTok{(}\AttributeTok{x =}\NormalTok{ mejoresB,}
                    \AttributeTok{y =}\NormalTok{ mejoresC,}
                    \AttributeTok{paired =} \ConstantTok{TRUE}\NormalTok{,}
                    \AttributeTok{alternative =} \StringTok{"two.sided"}\NormalTok{,}
                    \AttributeTok{conf.level =} \FloatTok{0.95}\NormalTok{,}
                    \AttributeTok{mu =} \DecValTok{0}
\NormalTok{                    )}
\end{Highlighting}
\end{Shaded}

\begin{verbatim}
## Warning in wilcox.test.default(x = mejoresB, y = mejoresC, paired = TRUE, :
## cannot compute exact p-value with ties
\end{verbatim}

\begin{Shaded}
\begin{Highlighting}[]
\NormalTok{prueba}\SpecialCharTok{$}\NormalTok{p.value}
\end{Highlighting}
\end{Shaded}

\begin{verbatim}
## [1] 0.002574545
\end{verbatim}

Dado que el valor p obtenido por la prueba es menor a nuestro alfa =
0.05, se rechaza la hipótesis nula en favor de la hipótesis alternativa,
por ende, las sospechas de la memorista son ciertas, es decir, las
mejores soluciones encontradas para la versión B y C del algoritmo
cuando estas poseen instancias de 70 o más nodos, poseen rendimientos
distintos.

\end{document}
